\indent 
    
    We seek the minimiser
\[
  x^* \;=\; \arg\min_{x \in \mathbb{R}}\; f(x),
\]
where $f$ is assumed differentiable and the derivative $f'$ is known analytically
(or can be computed).

\subsection{The update rule}

Starting from an initial guess $x_0$, gradient descent produces the sequence
\[
  \boxed{x_{k+1} \;=\; x_k \;-\; r \cdot f'(x_k)}, \qquad k = 0, 1, 2, \ldots
\]
where $r > 0$ is the \textbf{learning rate} (also called step size).

\paragraph{Intuition.}
The derivative $f'(x_k)$ points uphill; subtracting a fraction of it moves
$x$ downhill.  If $r$ is too small the progress is slow; if $r$ is too large
the iterates may overshoot and oscillate or diverge.

\subsection{Convergence criterion}

We stop when the gradient magnitude drops below a tolerance $\varepsilon$:
\[
  |f'(x_k)| \;<\; \varepsilon \;\approx\; 10^{-8}.
\]

\subsection{Test function}

Throughout this report we use:
\[
  f(x) = x^2, \qquad f'(x) = 2x.
\]
The unique global minimum is $x^* = 0$ with $f(x^*) = 0$.